\chapter{Seeds Embedded in Fruits}

Many plants produce their seeds embedded in a fruit. The natural cycle is for a
bird or animal to eat the fruit and to excrete the seed or to carry the fruit off and eat off
the pulp in a secluded spot. This is an effective mechanism for seed dispersal. It also
introduces a new mechanism for preventing germination of the seed before the seed is
dispersed. This-mechanism is to have the pulp contain powerful germination inhibitors
and to have a water channel through the seed coat to allow the inhibitor to
continuously diffuse into the embryo. When the pulp is removed, the embryo soon
destroys the remaining inhibitors and germination begins. Before going further, it is
important to point out that not all seeds in fruits use this mechanism, and the other
behaviors will be taken up later in this Chapter.

The presence of an inhibitor in the fruit is clearly the situation with Actinidia
chinensis (Kiwi fruit), Asparagus officinalis, Celastrus scandens (bittersweet),
Lycopersicon esculentum (the tomato), and Morus alba. Germination begins after
several days of washing with water using daily rinsing. From the viewpoint of chemical
kinetics these results shows that the inhibitor is continuously diffusing into the embryo,
that the embryo rapidly destroys the inhibitors by a chemical reaction, and that the half-
life of the inhibitors after diffusing into the seed must be less than one day.

A variation of the above is shown by Berberis thunbergii, Lonicera maacki and
sempervirens, and several Malus hybrids. After washing and cleaning for seven days,
the seed germinates at 40, but only after asignificant induction period indicating that
the half-life of the inhibitors in the seed is longer so that more time is needed for their
complete destruction. Incidentally these species do not germinate at 70.
All the evidence shows that the inhibition is due to specific chemical
compounds. The acidity can be neutralized and the osmotic pressure varied with no
effect on the inhibition (ref.14). The juice of fruits such as apple, grape,and lemon will
inhibit the germination of seeds that normally are immediate germinators (ref. 15).
Observant readers will have noticed that the seeds inside of citrus fruits and tomatoes
sometimes germinate if the fruit has been stored for a lengthy period. Rather than
being an exception these situations further support the picture presented above. After
lengthy periods the supply of inhibitor finally becomes exhausted and germination
begins.

Seeds of sweet peppers (Capsicum) and squash (Cucurbita) germinated 100\%
in 4-6 days after rinsing the seeds in water for just ten minutes. Not only was the
requisite washing period so short, but the seeds do not seem to be in contact with the
juices in the interior of the fruit. The seeds are attached by a thin filament. It is
possible that the germination inhibitors are continuously transmitted through this
filament. As soon as this filament is broken the seeds germinate.
The question arises as to whether simple water washing is enough or whether
surfactants and other chemicals in the digestive tract of birds and animals play a
positive role in germinating the seed. It was found that Euonymus alatus and
Euonymus europaeus germinated only if washed with aqueous detergents. so that the
inhibitor must be an oil soluble chemical compound. These two Euonymus have oily
fruits so that it is not surprising that the germination inhibitors are oil soluble.

The examples of seeds that germinate immediately after washing are the
exception. More often seeds from fruits have extended multicycle germinations. This
is found in Cornus, Cotoneaster, hex, Prunus, and Viburnum for example. With these
species the pulp may or may not have germination inhibitors since the seed already
has another mechanism for delaying germination until after seed is dispersed such as
the mechanisms described in Chapters 5-7. The question arises as to whether
grinding a hole through the seed coats would promote germination, particularly since
such seeds have hard and even bony seed coats and since the acid and enzymes in
the digestive tracts of birds and animals might weaken such seed coats, make them
more permeable, and promote germination. This subject of impervious seed coats is
treated in the following chapter, Chapter 9, but it is noted here that impervious seed
coats in seeds embedded in fruits are certainly rare. Prinsepia is the one clear
example of seed embedded in a fleshy fruit that has an impervious seed coat.
Generally treatment with gibberelic acid did not initiate germination in the seeds
from fruits. A notable exception was Viburnum trilobum. Its behavior is described in
Chapter 20.

All seeds embedded in fruits were washed and cleaned before the germination
experiments were started as explained in Chapter 3. This included species of
Cactaceae and Trillium. Although these are not usually thought of as fruits, I have
(cautiously) tasted a number of these fruits and found them to be sweet and of an
acceptable flavor. It is most probable that these depend on animals, birds, or insects
who eat the fruits and are responsible for dispersal of the seed. Claude Barr told me
that fox eat Opuntia fruits.

Some seeds have a fleshy appendage called an aril. Examples are seeds of
Chionodoxa, Jeffersonia, Trillium, and Oncocyclus and Regehia Iris. It is convenient to
remove the arils as this reduces the mold in thetowels, however, there is a more
fundamental question. The aril is designed to attract ants and other insects who carry
off the seed with the attached aril. They eat the aril at their leisure in a sort of insect
picnic. Two questions arise. Does the aril have to be removed for germination, and do
the ants and other insects secrete any chemicals that facilitate germination.
Experiments were performed to see whether removal of the aril affected germination.
Generally there was little if any effect. An exception was Chionbdoxa hucihiae where
germination occurred only with seed in which the aril had been soaked off.

Seed from botanical expeditions, seed exchanges, and commercial houses are
usually distributed in the dried fruit: Just as with fresh fruit the seed must be soaked,
washed, and cleaned. A few experiments were performed in which a comparison was
made between storing the seeds in the dried fruit relative to storing the seed in a
washed and cleaned state. Generally no differences were detected, and storing the
seed in the washed state gives less trouble from mold. Where there was a difference
as in Arisaema dracontium, storage in the dried fruit gave better storage life.
With all seeds in berries and fruits there is the possibility that the washing and
cleaning were not carried out for enough time. Some experiments were conducted
comparing washing and cleaning for seven days with washing and cleaning for four
weeks. The results are summarized in Table 8-1. The results can be complicated if
the germination inhibitors are oil soluble chemical species as in some Euonymus.

\begin{table}[h!]
\refstepcounter{table}
\begin{center}
Table \thetable: Comparison of Washing and Cleaning for Seven Days Versus Four Weeks
\end{center}

\begin{tabular}{|l|p{4cm}|p{4cm}|}
\hline
Species                & \multicolumn{2}{|c|}{Germination after washing and cleaning} \\
\hline
                       & Seven days   & Four weeks    \\
\hline
Physalis alkekengi (a) & 70L(43\%)    & 70L(33\%)     \\
                       & 40-70L(65\%) & 40-70L(68\%)  \\
Solanum dulcamara (a)  & 70L(93\%)    & 70L(97\%)     \\
                       & 70D none     & 70D none      \\
                       & 40-70L(100\%) & 40-70L(100\%)         \\
                       & 40-70D(23\%)  & 40-70D(60\%)          \\
Iris douglasiana DS 3 y& 70(1/3)-40(2/3) & 70(14/17)           \\
                       & 40(2/5)-70    & 40(13/13)             \\
Iris lactea DS 2 y     & 70-40-70-40   & 70-40-70(5\%)         \\
                       & 40-70-40-70(9\%) & 40-70(3\%)-40      \\
            & sown outdoors in Nov., \newline 60\% germ. in April-Sept. & sown outdoors in Nov., \newline  89\% germ. in April-May \\
Iris latifolia DS      & 70(10\%)-40(80\%) & 70-40(90\%)      \\
                       & 40-70-40(80\%)   & 40(50\%)-70-40(17\%)      \\
Iris magnifica DS 2 y  & 70-40(1/8)-70-40-70(1/8) & 70-40(15\%)      \\
                       & 40(93\%)         & 40(43\%)-70-40(17\%)      \\
                       & sown outdoors (30\%) & sown outdoors (30\%)      \\
Iris magnifica DS 9 y  & 70, all rotted   & 70-40(5\%)-70-40(5\%)      \\
                       & 40(100\%)     & 40-70(17\%)-40-70(4\%)      \\
                       & sown outdoors (15\%) & sown outdoors (50\%)      \\
Iris reticulate DS     & 70-40-70-40(2/5)     & 70-40-70-40(2/7)      \\
                       & 40-70-40-70      & 40-70-40-70      \\
Iris setosa DS 2 y     & 70(16\%)-40      & 70(12\%)-40      \\
                       & 40(4\%)-70(71\%) & 40(2\%)-70       \\
            & sown outdoors in Nov., \newline 97\% germ. in March & sown outdoors in Nov., \newline 87\% germ. in March and April     \\
Iris subbiflora DS 3 y & 70(2/12)-40-70  & 70(1/10)-40       \\
                       & 40(2/8)-70(1/8) & 40-70             \\
\hline
\end{tabular}

(a) The symbols L and D refer to light and
are described in Chapter 3. There was no germination of Physalis in dark.
\label{table:wash}
\end{table}

The reason for applying extended washing to so many Iris was that it had been
suggested that germination of Iris, particularly seed that had been dry stored for years
might benefit from extended soaking before sowing. In general the extended washing
had little effect except with Iris douglasiana and Iris tenax where some increase in
percent germination and rate of germination occurred. The extended soaking actually
reduced germination at 40 in Iris magnifica that had been dry stored for nine years.

Many Ericaceae have the seeds embedded in fruits, but the fruit does not
appear to contain inhibitors. These seeds depend on exposure to light to destroy the
inhibitor systems. Careful examination of the fruit of the common cranberry, Vaccinium
macrocarpum, will show that the seed is in a locule inside thEfruit. It can actually rattle
around inside this locule and has no direct contact with the juice of the fruit.

